\section{Organización del código y pipeline}
\label{sec:codebase}

Ambos trabajos siguen la misma convención de estructura usada en el
resto de \texttt{WORK/} del repositorio: un paquete Python por
laboratorio, con un módulo por responsabilidad en lugar de un único
archivo monolítico, un \texttt{main.py} que orquesta la ejecución de
las tareas numeradas y persiste las figuras resultantes, y una carpeta
\texttt{tests/} independiente con verificaciones numéricas reales
sobre propiedades de convergencia conocidas de antemano, en lugar de
pruebas que solo comprueban que el código no lance una excepción. Esta
separación entre el código de demostración (\texttt{main.py}, pensado
para producir las figuras de este informe) y el código de
verificación (\texttt{tests/}, pensado para ejecutarse de forma
automatizada y repetible) refleja una práctica deliberada de
ingeniería de software aplicada a experimentación numérica: el hecho
de que una figura "se vea razonable" no sustituye a una aserción
numérica explícita sobre el comportamiento esperado del algoritmo.

\begin{figure*}[!t]
\centering
\begin{verbatim}
WORK_ELEVEN/                              WORK_TWELVE/
|-- lms_lab/                              |-- nlms_lab/
|   |-- signals.py    (senales de prueba) |   |-- signals.py    (senales de prueba)
|   |-- unknown_system.py (sistema FIR    |   |-- noise.py      (ruido aditivo)
|   |            desconocido a identificar)   |-- adaptive.py   (NLMS y LMS, desde cero)
|   |-- lms.py        (regla LMS, desde   |   |-- audio.py      (carga perezosa librosa)
|   |            cero: run_lms)           |   `-- plotting.py
|   |-- audio.py      (carga perezosa     |-- audio_samples/gato/
|   |            librosa)                 |-- tests/
|   `-- plotting.py                       |   |-- test_adaptive.py
|-- audio_samples/gato/                   |   `-- test_signals.py
|-- tests/                                `-- main.py
|   |-- test_signals.py
|   |-- test_unknown_system.py
|   |-- test_lms.py
|   `-- test_mu_convergence.py
`-- main.py
\end{verbatim}
\caption{Estructura de directorios de WORK\_ELEVEN (LMS) y WORK\_TWELVE (NLMS). Ambas carpetas son completamente independientes entre sí: ninguna importa código de la otra.}
\label{fig:codebase_tree}
\end{figure*}

El flujo de procesamiento (\textit{pipeline}) es el mismo en ambos
trabajos y se ejecuta íntegramente dentro de \texttt{main.py} al
invocar cada tarea:

\begin{enumerate}
  \item \textbf{Generación de señales} (\texttt{signals.py}): se
  construyen las señales de prueba (tono, ruido blanco, señal de
  potencia variable) y, para las tareas de cancelación de ruido, se
  carga el audio real de referencia (\texttt{audio.py}).
  \item \textbf{Sistema/ruido de referencia}: en el caso de
  identificación de sistema, \texttt{unknown\_system.py} (WORK\_ELEVEN)
  o la definición equivalente en \texttt{adaptive.py} (WORK\_TWELVE)
  aplica un filtro FIR desconocido a la señal de entrada para producir
  la señal deseada $d[n]$; en el caso de cancelación de ruido,
  \texttt{noise.py} agrega ruido aditivo correlacionado a la señal
  limpia.
  \item \textbf{Filtrado adaptativo}: \texttt{lms.py}
  (\texttt{run\_lms}, WORK\_ELEVEN) o \texttt{adaptive.py}
  (\texttt{nlms\_filter}, WORK\_TWELVE) ejecuta la regla de
  actualización correspondiente (Ec.~\ref{eq:lms} o
  Ec.~\ref{eq:nlms}) muestra a muestra, devolviendo la salida
  estimada, el error en cada iteración y la evolución de los pesos.
  \item \textbf{Visualización} (\texttt{plotting.py}): cada tarea de
  \texttt{main.py} llama a una función de graficado dedicada que
  guarda la figura correspondiente en \texttt{output/}.
  \item \textbf{Verificación} (\texttt{tests/}): de forma
  independiente al pipeline de \texttt{main.py}, los tests de
  \texttt{pytest} ejecutan el mismo algoritmo sobre casos sintéticos
  con resultado conocido de antemano (p.\,ej. coeficientes de un
  sistema FIR fijo) y comprueban numéricamente la convergencia, en
  lugar de solo verificar que el código no lance una excepción.
\end{enumerate}

En WORK\_ELEVEN, \texttt{lms\_lab/lms.py} concentra la totalidad de la
lógica del algoritmo en una sola función (\texttt{run\_lms}), que
recibe la señal de entrada, la señal deseada y el tamaño de paso, y
devuelve tanto la historia completa de pesos como la señal de error —
exponiendo deliberadamente el estado interno del filtro en cada
iteración, en lugar de devolver únicamente el resultado final, para
que las funciones de graficado y las pruebas puedan inspeccionar la
trayectoria completa de convergencia y no solo su punto final.
\texttt{unknown\_system.py} define los coeficientes del sistema FIR
desconocido como una constante del módulo, de forma que tanto
\texttt{main.py} como los tests de \texttt{tests/test\_unknown\_system.py}
referencian el mismo valor de verdad conocida, evitando duplicar un
número "mágico" en dos lugares que podrían desincronizarse.

En WORK\_TWELVE, \texttt{nlms\_lab/adaptive.py} contiene dos funciones:
\texttt{nlms\_filter}, con la regla de actualización normalizada
(Ec.~\ref{eq:nlms}), y \texttt{lms\_filter}, una reimplementación
mínima del algoritmo LMS plano usada únicamente como punto de
comparación en la Sección~\ref{sec:nlms_vs_lms} —se decidió duplicar
esta pequeña pieza de código en lugar de importarla de WORK\_ELEVEN
precisamente para que cada carpeta de trabajo permanezca
autocontenida y pueda ejecutarse, probarse y calificarse de forma
completamente independiente, sin que un cambio futuro en un trabajo
pueda romper silenciosamente al otro. \texttt{noise.py} centraliza la
generación de ruido aditivo gaussiano reutilizada por las tareas de
cancelación de ruido y por los propios tests.
